\section{Related Work}
\label{sec:related}

\paragraph{Video-generation evaluation.}
Early distributional measures such as Fr\'echet Video Distance compare generated and real feature distributions~\cite{unterthiner2018fvd}, but a single corpus-level statistic offers limited diagnostic value for prompt-conditioned generation. VBench decomposes generation quality into multiple perceptual and semantic dimensions and validates them against human preference~\cite{huang2024vbench}; VBench++ extends this framework to image-to-video settings and trustworthiness~\cite{huang2025vbenchpp}. EvalCrafter combines a broad prompt set with objective metrics calibrated to human opinion~\cite{liu2023evalcrafter}. In contrast, \FORGE focuses on industrial correctness: structural invariants, event causality, physical evolution, camera execution, and downstream decision value.

\paragraph{Compositional and reasoning-oriented evaluation.}
T2V-CompBench tests attribute, action, interaction, spatial, and numerical composition using learned and programmatic metrics~\cite{sun2024t2vcompbench}. This direction motivates evaluating more than generic text--video similarity. Industrial prompts, however, contain constraints that are often conditional: a machine should stop \emph{after} a boundary is crossed, a fracture should remain localized, or a load should move while its rigging stays coupled. \FORGE therefore associates each task with observable events, implicit-rule questions, failure modes, and task-specific evidence.

\paragraph{Multimodal judges and hybrid metrics.}
Large multimodal models can describe temporal evidence and apply natural-language rubrics, but their outputs may be sensitive to prompt wording, incomplete visual sampling, and score anchoring. Purely programmatic metrics have the opposite limitation: they are reproducible but capture only a narrow proxy. Our protocol combines the two. A model judge assigns the semantic axis scores, while classical operators expose measurements such as joint drift, regional change, plume continuity, and camera motion. Operator evidence is validity- and confidence-gated and is never counted as an additional independent axis.

\paragraph{Industrial visual understanding.}
Industrial perception research commonly targets a single downstream problem, such as defect inspection, robot manipulation, safety monitoring, or predictive maintenance. \FORGE instead evaluates a generator's ability to produce videos that could plausibly support these workflows. It does not claim that a generated clip is safe for deployment; rather, it measures whether the visual evidence satisfies a frozen industrial task contract and exposes where that contract fails.

\paragraph{Comparison with physics-oriented benchmarks.}
Table~\ref{tab:benchmark-comparison} compares representative physics-oriented video benchmarks by scenario, data source, physical depth, topology-failure coverage, evaluation pipeline, and scoring mechanism. \FORGE differs in targeting high-risk industrial scenarios, explicitly testing topology failures, and combining objective computer-vision operators with multimodal causal question answering under a strict static gate.

\begin{table*}[t]
\centering
\scriptsize
\setlength{\tabcolsep}{2.4pt}
\renewcommand{\arraystretch}{1.12}
\caption{Comparison of representative physics-oriented video benchmarks. Coverage is described as full, partial, or none without symbol-dependent typography.}
\label{tab:benchmark-comparison}
\resizebox{\textwidth}{!}{%
\begin{tabular}{lllllll}
\toprule
Benchmark & Scenario domain & Data source & Physics depth & Topology failure & Evaluation pipeline & Core metric \\
\midrule
PhyGenBench (2025) & High-school physics / common sense & Web images / synthetic & Partial: basic laws & None & VLM visual QA & Weighted mean \\
RigidBench (2026) & Synthetic rigid-body collisions & Blender 3D rendering & Full: rigid trajectories & None & 3D trajectory error & MAE \\
Physion-Eval (2026) & Daily life / human actions & Real-world short videos & Full: daily physics & None & VLM failure detection & Failure rate / aggregate \\
Physics-IQ (2025) & General video dynamics & Generic video clips & Partial: macro dynamics & None & VLM multiple choice & Accuracy \\
VBench-2.0 (2025) & General generation & Generic generations & Partial: shallow intuition & None & Pretrained VLM judge & Mean across dimensions \\
\textbf{\FORGE{} (ours)} & \textbf{Industrial / security / disaster / manufacturing} & \textbf{Surveillance / satellite / BIM / scientific diagrams} & \textbf{Full: extreme dynamics / fluids} & \textbf{Full: targeted} & \textbf{CV operators + VLM causal QA} & \textbf{Strict usability with static veto} \\
\bottomrule
\end{tabular}%
}
\end{table*}
